\section{AIME II 2026}
        \begin{enumerate}
        	\item (\textit{AIME II 2026}) Find the sum of the $10$th terms of all arithmetic sequences of integers that have first term equal to $4$ and include both $24$ and $34$ as terms.


 

	\item (\textit{AIME II 2026}) The figure below shows a grid of $10$ squares in a row. Each square has a diagonal connecting its lower left vertex to its upper right vertex. A bug moves along the line segments from vertex to vertex, never traversing the same segment twice and never moving from right to left along a horizontal or diagonal segment. Let $N$ be the number of paths the bug can take from the lower left corner ($A$) to the upper right corner ($B$). One such path from $A$ to $B$ is shown by the thick line segments in the figure. Find $\sqrt{N}$.


 
                    \begin{figure}[H]
                        \centering
                        \includegraphics[width=0.5\textwidth]{Problem/AIME_II_2026_2_Problem_Figure_1.png}
                    \end{figure}
                    


	\item (\textit{AIME II 2026}) Let $ABCDE$ be a nonconvex pentagon with internal angles $\angle A = \angle E = 90^\circ$ and $\angle B = \angle D = 45^\circ.$ Suppose that $DE < AB, AE = 20, BC = 14\sqrt2,$ and points $B,C,$ and $D$ lie on the same side of line $AE.$ Suppose further that $AB$ is an integer with $AB < 2026$ and the area of pentagon $ABCDE$ is an integer multiple of $16.$ Find the number of possible values of $AB.$


	\item (\textit{AIME II 2026}) For each positive integer $n$ let $f(n)$ be the value of the base-ten numeral $n$ viewed in base $b$, where $b$ is the least integer greater than the greatest digit in $n$. For example, if $n=72$, then $b=8$, and $72$ as a numeral in base $8$ equals $7\cdot 8+2=58$; therefore $f(72)=58$. Find the number of positive integers $n$ less than $1000$ such that $f(n)=n$.


 

	\item (\textit{AIME II 2026}) An urn contains $n$ marbles. Each marble is either red or blue, and there are at least $7$ marbles of each color. When $7$ marbles are drawn randomly from the urn without replacement, the probability that exactly $4$ of them are red equals the probability that exactly $5$ of them are red. Find the sum of the five least values of $n$ for which this is possible.


 

	\item (\textit{AIME II 2026}) Find the sum of all real numbers $r$ such that there is at least one point where the circle with radius $r$ centered at $(4, 39)$ is tangent to the parabola with equation $2y = x^2 - 8x + 12.$


	\item (\textit{AIME II 2026}) A standard fair six-sided die is rolled repeatedly. Each time the die reads 1 or 2, Alice gets a coin; each time it reads 3 or 4, Bob gets a coin; and each time it reads 5 or 6, Carol gets a coin. The probability that Alice and Bob each receive at least two coins before Carol receives any coins can be written as $\frac mn,$ where $m$ and $n$ are relatively prime positive integers. Find $100m+n.$


	\item (\textit{AIME II 2026}) Isosceles triangle ${}\triangle ABC$ has $AB = BC.$ Let $I$ be the incenter of ${}\triangle ABC.$ The perimeters of ${}\triangle ABC$ and ${}\triangle AIC$ are in the ratio $125:6,$ and all the sides of both triangles have integer lengths. Find the minimum possible value of $AB.$


	\item (\textit{AIME II 2026}) Let $S$ denote the value of the infinite sum
 \[\frac 19 + \frac 1{99} + \frac 1{999} + \frac 1{9999}  + \cdots\]Find the remainder when the greatest integer less than or equal to $10^{100} S$ is divided by $1000.$


	\item (\textit{AIME II 2026}) Let $\triangle ABC{}$ be a triangle with $D$ on $\overline{BC}$ such that $\overline{AD}{}$ bisects $\angle BAC.$ Let $\omega{}$ be the circle that passes through $A$ and is tangent to segment $\overline{BC}{}$ at $D.$ Let $E \neq A$ and $F \neq A$ be the intersections of $\omega{}$ with segments $\overline{AB}{}$ and $\overline{AC},$ respectively. Suppose that $AB = 200, AC = 225,$ and all of $AE, AF, BD,$ and $CD{}$ are positive integers. Find the greatest possible value of $BC.$


	\item (\textit{AIME II 2026}) Find the greatest integer $n$ such that the cubic polynomial
 \[x^{3} -\frac{n}{6}x^{2} + (n - 11)x - 400\]has roots $\alpha^{2}$, $\beta^{2}$, and $\gamma^{2}$, where $\alpha$, $\beta$, and $\gamma$ are complex numbers, and there are exactly seven different possible values for $\alpha + \beta + \gamma$.


 

	\item (\textit{AIME II 2026}) Consider a tetrahedron with two isosceles triangle faces with side lengths $5\sqrt{10}, \,5\sqrt{10},$ and $10$ and two isosceles triangle faces with side lengths $5\sqrt{10},\, 5\sqrt{10},$ and $18.$ The four vertices of the tetrahedron lie on a sphere with center $S,$ and the four faces of the tetrahedron are tangent to a sphere with center $R.$ The distance $RS$ can be written as $\frac mn,$ where $m$ and $n$ are relatively prime positive integers. Find ${}m + n.$


	\item (\textit{AIME II 2026}) Call finite sets of integers $S$ and $T$ cousins if


 
            \begin{itemize}
        \item $S$ and $T$ have the same number of elements,

        \item $S$ and $T$ are disjoint, and

        \item the elements of $S$ can be paired with the elements of $T$ so that the elements in each pair differ by exactly $1$.
            \end{itemize}
            For example, $\{1,2,5\}$ and $\{0,3,4\}$ are cousins. Suppose that the set $S$ has exactly $4040$ cousins. Find the least number of elements the set $S$ can have.


 

	\item (\textit{AIME II 2026}) For integers $a$ and $b,$ let $a \circ b = a - b$ if $a$ is odd and $b$ is even, and $a+b$ otherwise. Find the number of sequences $a_1, a_2, a_3, \ldots, a_n$ of positive integers such that
 \[a_1 + a_2 + a_3 + \cdots + a_n = 12 \quad \text{and} \quad a_1 \circ a_2 \circ a_3 \circ \cdots \circ a_n = 0\]where the operations are performed from left to right; that is, $a_1 \circ a_2 \circ a_3$ means $(a_1 \circ a_2) \circ a_3.$


	\item (\textit{AIME II 2026}) Find the number of ordered 7-tuples $(a_1, a_2, a_3, \ldots, a_7)$ having the following properties:


 
            \begin{itemize}
        \item $a_k \in \{1,2,3\}$ for all $k.$

        \item $a_1+a_2+a_3+a_4+a_5+a_6+a_7$ is a multiple of $3.$

        \item $a_1a_2 a_4 + a_2a_3a_5 + a_3a_4 a_6 + a_4 a_5 a_7 + a_5 a_6 a_1 + a_6 a_7 a_2 + a_7 a_1 a_3$ is a multiple of $3.$
            \end{itemize}
            

\end{enumerate}